% !TeX root = ../../main.tex
\section{Relazione tra Misura, Variabile Aleatoria e Distribuzione}

Per fare chiarezza su quanto detto, analizziamo le differenze e le relazioni tra i concetti chiave.

\paragraph{Definizioni Fondamentali}
\begin{itemize}
    \item \textbf{La Misura di Probabilità ($\mathbb{P}$):} È una funzione che prende in ingresso un \textit{evento} (un insieme di risultati) e restituisce un numero reale compreso tra 0 e 1. Essa quantifica quanto è probabile che accada quel certo evento.
    
    \item \textbf{La Variabile Aleatoria ($X$):} È una funzione che prende un singolo risultato astratto ($\omega \in \Omega$) e lo trasforma in un numero reale.
    \\Serve a tradurre un risultato "fisico" o astratto in un valore quantitativo.
    \\\textit{Esempio:} "Esce la faccia del Re di cuori" $\xrightarrow{X}$ "vale 10 punti".
\end{itemize}

\subsection{La Distribuzione (Misura Indotta)}

La misura di probabilità indotta sullo spazio misurabile di arrivo $(E, \mathcal{F})$ da una variabile aleatoria $X$, a partire dalla misura originale $\mathbb{P}$ su $(\Omega, \mathcal{E})$, è detta \textbf{distribuzione}, \textbf{legge} o \textbf{probabilità di X}, ed è indicata con $P_X$.

In pratica, è la misura associata alla variabile aleatoria. Mentre la misura originale $\mathbb{P}$ agisce sugli eventi dello spazio campionario, la distribuzione $P_X$ restituisce la probabilità dei valori numerici in output.

\textbf{Concetto chiave:}
\begin{quote}
    È la misura $\mathbb{P}$ che è stata \textbf{"trasportata"} sui numeri reali tramite la variabile aleatoria $X$.
\end{quote}


Ovviamente, la forma e il comportamento della distribuzione $P_X$ dipendono interamente dalla definizione della variabile aleatoria $X$.

\subsection{Tipologie di Distribuzioni}

Le distribuzioni si dividono in due grandi famiglie in base alla natura del loro codominio:

\begin{enumerate}
    \item \textbf{Discrete:} Se la funzione $X$ mappa i risultati su un insieme di numeri che si possono \textit{contare} (insieme numerabile).
    \item \textbf{Continue:} Se il codominio è un intervallo continuo di numeri reali.
\end{enumerate}

A differenza delle variabili continue (dove la probabilità in un singolo punto esatto è matematicamente zero), nelle distribuzioni discrete la probabilità è concentrata in \textbf{"grumi"} o \textbf{"masse"} su specifici punti.

La funzione di massa di probabilità $p_X(x)$ risponde alla domanda:
\[
\text{"\textbf{Esattamente quanta probabilità è accumulata su questo specifico numero?}"}
\]

\subsection{La Variabile di Bernoulli}

\paragraph{Definizione e Spazio Campione}
Consideriamo un esperimento che prevede una sola esecuzione ($N = 1$), come il singolo lancio di una moneta. Lo spazio campione in questo caso si riduce agli esiti di base: $\Omega = \{C, T\}$.
\\Definiamo la variabile aleatoria indicatrice $X$, la quale mappa l'esito dell'esperimento in un dominio numerico strettamente binario:
\[ X: \omega \in \Omega \longrightarrow X(\omega) \in \{0, 1\} \]
Convenzionalmente, associamo il valore $1$ al successo ("Testa") e il valore $0$ all'insuccesso ("Croce").

\paragraph{Funzione di Massa di Probabilità (PMF)}
La distribuzione di probabilità per una variabile di Bernoulli, denotata con la scrittura $X \sim \mathcal{B}(p)$, è definita in modo assiomatico dalle probabilità dei due unici eventi possibili:
\[ 
\mathbb{P}(X = x) = 
\begin{cases} 
p & \text{se } x = 1 \\ 
q = 1 - p & \text{se } x = 0 
\end{cases} 
\]
Al fine di operare analiticamente, possiamo unificare i due casi in un'unica espressione compatta, valida per l'intero alfabeto $x \in \{0, 1\}$:
\[ p_X(x) = p^x q^{1-x} \]

\subsection{Esempio: Il Conteggio Bernoulliano}

Per applicare concretamente i concetti appena esposti, consideriamo un esperimento fondamentale nella teoria della probabilità: il lancio ripetuto di una moneta. Questo scenario è il prototipo del cosiddetto \textbf{Processo di Bernoulli}.

\paragraph{Descrizione dell'Esperimento}
Immaginiamo di lanciare una moneta per $N$ volte consecutive. Possiamo modellare questo fenomeno attraverso tre ipotesi chiave:
\begin{itemize}
    \item \textbf{Esiti binari:} Ogni singolo lancio ammette solo due risultati: "Testa" ($T$), che chiameremo successo, e "Croce" ($C$), che chiameremo insuccesso.
    \item \textbf{Probabilità costanti:} La moneta ha una probabilità $p$ (con $0 < p < 1$) di dare "Testa". Di conseguenza, la probabilità di ottenere "Croce" sarà $q = 1 - p$.
    \item \textbf{Indipendenza:} Assumiamo che i lanci siano statisticamente indipendenti; l'esito di un lancio non influenza in alcun modo quelli precedenti o successivi.
\end{itemize}

\paragraph{Ora Passiamo dallo Spazio Campione alla Variabile Aleatoria}
Lo spazio campione $\Omega$ è costituito dall'insieme di tutte le possibili sequenze di risultati (le $N$-ple). Un evento elementare generico $\omega \in \Omega$ si presenta come una stringa ordinata del tipo:
\[ \omega = (C, C, T, T, C, \dots, C, T, T) \]

Il nostro obiettivo non è analizzare l'ordine esatto, ma il numero complessivo di successi. Definiamo quindi la variabile aleatoria $X_N$ che "conta" quante volte esce "Testa" negli $N$ lanci:
\[ X_N: \omega \in \Omega \longrightarrow X_N(\omega) \in \{0, 1, \dots, N\} \]
Diremo che $X_N(\omega) = k$ se la sequenza $\omega$ contiene esattamente $k$ successi.


\subsubsection{Costruzione della Probabilità}

Per determinare la probabilità che si verifichino esattamente $k$ successi ($P(X_N=k)$), procediamo per gradi sfruttando l'indipendenza degli eventi.

\paragraph{1. Analisi di una singola configurazione}
Consideriamo una specifica sequenza $\omega^*$ che soddisfa la condizione di avere $k$ teste e $N-k$ croci. Per fissare le idee, immaginiamo la sequenza in cui appaiono prima tutte le teste e poi tutte le croci. Grazie all'indipendenza stocastica, la probabilità congiunta è semplicemente il prodotto delle singole probabilità:

\[
\mathbb{P}(\omega^*) = \mathbb{P}(\underbrace{T \cap \dots \cap T}_{k} \cap \underbrace{C \cap \dots \cap C}_{N-k}) = p^k q^{N-k}
\]
Questa rappresenta la probabilità di una singola, specifica sequenza in cui i risultati si presentano in un ordine ben preciso.
È fondamentale notare che qualsiasi altra sequenza con lo stesso numero di teste e croci (anche se in ordine diverso) avrà la \textbf{stessa identica probabilità}, poiché il prodotto è commutativo.

\paragraph{2. Il ruolo della Combinatoria}
Poiché gli eventi elementari sono mutuamente esclusivi, la probabilità totale dell'evento "ottenere $k$ teste" è la somma delle probabilità di tutte le sequenze favorevoli.
\\Dobbiamo quindi chiederci: \textit{quante sono le sequenze distinte di lunghezza $N$ che contengono esattamente $k$ teste?}
\\La risposta è fornita dal \textbf{coefficiente binomiale}:
\[ C_{N,k} = \binom{N}{k} \]
L'insieme di queste sequenze costituisce l'evento composto $\Omega_{N,k}$.

\subsection{La Distribuzione Binomiale}

Unendo il conteggio combinatorio con la probabilità della singola sequenza, otteniamo la \textbf{Funzione di Massa di Probabilità} (PMF) Binomiale.
\\La probabilità che la variabile aleatoria $X_N$ assuma valore $k$ è data da:

\begin{equation}
\boxed{p_{X_N}(k) = \mathbb{P}(X_N = k) = \binom{N}{k} p^k q^{N-k}}
\end{equation}

Quando una variabile segue questa legge, si usa la notazione sintetica $X_N \sim \mathcal{B}(N, p)$.

\paragraph{Verifica di Normalizzazione}
Una proprietà essenziale di ogni PMF è che la somma delle probabilità su tutto l'alfabeto sia pari a 1. Nel nostro caso, sfruttando il celebre \textbf{Teorema del Binomio di Newton}, osserviamo che:

\[
\sum_{k=0}^{N} p_{X_N}(k) = \sum_{k=0}^{N} \binom{N}{k} p^k q^{N-k} = (p + q)^N
\]

Poiché $p$ e $q$ sono complementari ($p+q=1$), otteniamo $(1)^N = 1$, confermando che la distribuzione è ben definita.

\paragraph{Valore Atteso}
Infine, calcoliamo la media statistica della distribuzione. Applicando la definizione di valore atteso, si dimostra che:

\begin{equation}
\mathbb{E}[X_N] = \sum_{k=0}^{N} k \cdot \binom{N}{k} p^k q^{N-k} = Np
\end{equation}

Il risultato $\mathbb{E}[X_N] = Np$ è coerente con l'intuizione: se ripetiamo un esperimento $N$ volte con probabilità di successo $p$, ci aspettiamo in media $N \cdot p$ successi.
\\\\Alla luce di questa formulazione, risulta evidente come la Variabile Aleatoria Binomiale $X_N$, definita per $N$ lanci, non sia altro che la somma di $N$ variabili di Bernoulli indipendenti e identicamente distribuite (i.i.d.).

\subsection{La Distribuzione Uniforme Discreta}

Se la variabile di Bernoulli modellava una situazione con esattamente 2 esiti, la \textbf{Variabile Uniforme} modella situazioni in cui tutti gli esiti possibili hanno la \textbf{stessa probabilità} di verificarsi (sono equiprobabili).
\\Questa distribuzione concerne tutti gli esempi classici visti fino ad ora, come il lancio di un dado onesto, l'estrazione di una carta da un mazzo completo, ecc.

\paragraph{Definizioni}
Una variabile aleatoria $X$ si dice \textit{uniformemente distribuita} su un insieme (alfabeto) $\mathcal{X}$ (e si indica con $X \sim \mathcal{U}(\mathcal{X})$) se:
\begin{itemize}
    \item L'insieme dei valori possibili ha cardinalità finita $M$ (ossia $|\mathcal{X}| = M$).
    \item Ogni singolo valore $x \in \mathcal{X}$ ha probabilità $1/M$.
\end{itemize}

La sua funzione di massa di probabilità è quindi definita come:

\begin{equation}
\boxed{p_X(x) = \mathbb{P}(X=x) = \frac{1}{M}, \quad \forall x \in \mathcal{X}}
\end{equation}

Ovviamente $p_X(x)$ soddisfa le condizioni necessarie per essere una pmf (valori positivi e somma unitaria).

\paragraph{Esempio Classico}
Consideriamo il lancio di un dado a 6 facce onesto.
\begin{itemize}
    \item \textbf{Alfabeto:} $\mathcal{X} = \{1, 2, 3, 4, 5, 6\}$.
    \item \textbf{Cardinalità:} $M = 6$.
    \item \textbf{Probabilità puntuale:} La probabilità che esca un numero specifico (es. il 4) è:
    \[ \mathbb{P}(X=4) = \frac{1}{6} \]
\end{itemize}

\subsubsection{Valore Atteso (Media)}

Il calcolo della media è immediato. Applicando la definizione generale:

\begin{equation}
\mathbb{E}[X] = \sum_{x \in \mathcal{X}} x \cdot p_X(x) = \sum_{x \in \mathcal{X}} x \cdot \frac{1}{M} = \frac{1}{M} \sum_{x \in \mathcal{X}} x
\end{equation}

Poiché la probabilità $1/M$ è costante, la portiamo fuori dalla sommatoria. È ovvio notare che, in questo caso specifico, la media della variabile aleatoria \textbf{coincide con la media aritmetica} di tutti i possibili valori dell'alfabeto.

\subsubsection{Un Caso Notevole: Interi consecutivi da 0}

Un caso molto interessante, frequente in informatica, si ha quando l'alfabeto è costituito dai numeri interi da $0$ a $M-1$.
\[ \mathcal{X} = \{0, 1, \dots, M-1\} \]

Che succede al valore atteso in questo caso? Per calcolarlo sfruttiamo la \textbf{formula della somma dei primi interi} (sommatoria di Gauss).

\begin{enumerate}
    \item Dobbiamo sommare i numeri da $0$ a $M-1$.
    \item Ricordiamo che la somma da $1$ a $n$ è data da $\frac{n(n+1)}{2}$.
    \item Nel nostro caso il numero più alto è $n = M-1$.
    \item Sostituendo nella formula della somma:
    \[ \text{Somma} = \sum_{i=0}^{M-1} i = \frac{(M-1)((M-1)+1)}{2} = \frac{(M-1)M}{2} \]
    \item Ora calcoliamo la media dividendo per $M$ (come richiesto dalla formula della media uniforme):
    
    \begin{equation}
    \begin{aligned}
    \mathbb{E}[X] &= \frac{1}{M} \cdot \text{Somma} \\
    &= \frac{1}{M} \cdot \frac{M(M-1)}{2}
    \end{aligned}
    \end{equation}
    
    \item Semplificando $M$, otteniamo la formula finale per questo caso specifico:
    
    \begin{equation}
    \boxed{\mathbb{E}[X] = \frac{M-1}{2}}
    \end{equation}
\end{enumerate}

\paragraph{Esempio numerico}
Ipotizziamo di avere una cifra decimale generata a caso (da 0 a 9). In questo caso $M=10$.
\begin{itemize}
    \item I valori sono $\{0, 1, \dots, 9\}$.
    \item Applicando la formula appena ricavata:
    \[ \mathbb{E}[X] = \frac{10-1}{2} = \frac{9}{2} = 4.5 \]
    \item Il valore atteso è $4.5$.
\end{itemize}

\subsection{La Distribuzione di Poisson}

Prima di definire formalmente questa distribuzione, consideriamo uno scenario intuitivo per comprenderne l'utilità.

\paragraph{Esempio: Il Call Center}
Immaginiamo di lavorare in un call center. Analizzando i dati storici, sappiamo che, \textbf{in media}, riceviamo \textbf{5 chiamate all'ora}.
\\Tuttavia, il fatto che la media sia 5 non garantisce che nella prossima ora riceveremo \textit{esattamente} 5 chiamate. Essendo un fenomeno aleatorio:
\begin{itemize}
    \item In un'ora potremmo riceverne 3 (bassa attività).
    \item In un'altra potremmo riceverne 8 (picco di attività).
    \item In un'altra ancora potremmo non riceverne nessuna (0).
\end{itemize}

La distribuzione di Poisson nasce proprio per rispondere a domande specifiche su questi conteggi, del tipo:
\begin{quote}
    \textit{"Sapendo che la media è 5 chiamate/ora, qual è la probabilità che nella prossima ora ne arrivino esattamente 7?"}
\end{quote}

È la distribuzione di riferimento per modellare il conteggio di \textbf{eventi rari} che si verificano in un intervallo continuo di tempo, spazio o volume.

\paragraph{Il Parametro $\lambda$ e le Condizioni}
Tutta la distribuzione ruota attorno a un solo parametro fondamentale, $\lambda$ (Lambda), che rappresenta l'\textbf{intensità media} del processo:

\[ \lambda = \text{numero medio di eventi attesi nell'intervallo} \]

Affinché si possa applicare il modello di Poisson, devono valere tre ipotesi:
\begin{enumerate}
    \item \textbf{Indipendenza:} L'accadimento di un evento non influenza la probabilità che ne accada un altro (es. ricevere una chiamata ora non cambia la probabilità di riceverne una tra un minuto).
    \item \textbf{Stazionarietà:} Il tasso medio $\lambda$ è costante nel tempo considerato.
    \item \textbf{Non contemporaneità:} La probabilità che due eventi accadano nello stesso identico istante è trascurabile.
\end{enumerate}

\subsubsection{Derivazione come limite della Binomiale}

La formula di Poisson non è arbitraria, ma si ottiene matematicamente come caso limite della distribuzione Binomiale.

Immaginiamo di suddividere il nostro intervallo di tempo (l'ora del call center) in $N$ sotto-intervalli piccolissimi (es. millisecondi).
\begin{itemize}
    \item Sia $N \to \infty$ (numero di sotto-intervalli tendente all'infinito).
    \item Sia $p \to 0$ (probabilità di ricevere una chiamata nel singolo millisecondo è quasi nulla).
    \item Manteniamo però costante il prodotto $N \cdot p = \lambda$ (la media totale delle chiamate nell'ora resta 5).
\end{itemize}

Partiamo dalla PMF Binomiale e sostituiamo $p$ con $\frac{\lambda}{N}$:
N.B: ${p} = \frac{\lambda}{N}$
\[
\mathbb{P}(X=k) = \lim_{N \to \infty} \binom{N}{k} \left( \frac{\lambda}{N} \right)^k \left( 1 - \frac{\lambda}{N} \right)^{N-k}
\]

Espandiamo i termini per calcolare il limite:

\begin{equation}
\begin{aligned}
\mathbb{P}(X=k) &= \lim_{N \to \infty} \frac{N!}{k!(N-k)!} \cdot \frac{\lambda^k}{N^k} \cdot \left( 1 - \frac{\lambda}{N} \right)^N \cdot \left( 1 - \frac{\lambda}{N} \right)^{-k} \\
&= \frac{\lambda^k}{k!} \lim_{N \to \infty} \underbrace{\frac{N(N-1)\dots(N-k+1)}{N^k}}_{\to 1} \cdot \underbrace{\left( 1 - \frac{\lambda}{N} \right)^N}_{\to e^{-\lambda}} \cdot \underbrace{\left( 1 - \frac{\lambda}{N} \right)^{-k}}_{\to 1}
\end{aligned}
\end{equation}

Dove abbiamo sfruttato il limite notevole $\lim_{x \to \infty} (1 - \frac{\lambda}{x})^x = e^{-\lambda}$.

\paragraph{Funzione di Massa di Probabilità}
Il risultato del limite ci fornisce la formula della PMF per una variabile Poissoniana $X \sim \text{Pois}(\lambda)$:

\begin{equation}
\boxed{p_X(k) = \mathbb{P}(X=k) = \frac{\lambda^k \cdot e^{-\lambda}}{k!}, \quad k \in \{0, 1, 2, \dots\}}
\end{equation}

Dove $e$ è il numero di Nepero ($e \approx 2.718$).

\paragraph{Proprietà: Media e Varianza}
Una caratteristica unica ed elegante della distribuzione di Poisson è che il valore atteso e la varianza coincidono e sono pari al parametro stesso.

\begin{equation}
\mathbb{E}[X] = \lambda \quad \text{e} \quad \text{Var}(X) = \lambda
\end{equation}

Questo implica che in un processo di Poisson, se aumenta il numero medio di eventi attesi (es. un call center più grande), aumenta proporzionalmente anche l'incertezza (la variabilità) attorno a tale media.

